% Blueprint 内容：分解式节点（行内标准）。
%   \lean{声明}  链接到具体 Lean 声明（checkdecls 会校验存在）。
%   \leanok      仅用于「已被 Lean 内核完整证明、无 sorry」的节点 → 依赖图中显示为绿。
%                含 sorry / 库缺口的节点只给 \lean 链接、不给 \leanok，并在正文明写状态。
%   \uses{...}   标依赖，生成依赖图的边。
% 命题与证明的逐字原文见 docs/paper-lemma1.md（唯一权威来源）；本章为其形式化分解。

\chapter{Lemma 1: Posterior process}

本章形式化论文的 \textbf{Lemma 1 (Posterior process)}（论文标签
\texttt{lem:posterior-martingale}）。\textbf{命题与证明的逐字原文见 \texttt{docs/paper-lemma1.md}}；
本章把该证明**按论文自身的四步分解成独立节点**，每个节点链接到它对应的 Lean 声明，
并按是否被 Lean 内核完整证明标注状态（绿色 = 无 \texttt{sorry}，已背书）。读者可点开任一节点核对它证的是什么、以及 Lean 有没有真的证完。

\begin{definition}[Environment]\label{def:setup}
    离散时间 $t\in\mathbb T=\{0,1,2,\dots\}$。Borel 概率空间 $(\Omega,\mathcal F,\mathbb P)$，
    原始状态 $\theta:\Omega\to\Theta$，$\Theta$ 紧 Polish；$\Delta(\Theta)$ 为 $\Theta$ 上概率测度
    （弱拓扑）。公开状态 $Z_t:\Omega\to E_t$（$E_t$ Polish），$\mathcal F_t=\sigma(Z_t)$ 为递增滤子，
    $\mathcal F_0$ 平凡。
\end{definition}

%======================================================================
% 主命题（逐字原文），其证明拆为下面四个步骤节点。
%======================================================================
\begin{lemma}[Posterior process]\label{lem:posterior-martingale}
    \uses{def:setup}
    For each $t\in\mathbb T$, there exists an $\mathcal F_t$-measurable posterior kernel
    $S_t:\Omega\to\Delta(\Theta)$ such that, for every bounded Borel function
    $\varphi:\Theta\to\mathbb R$,
    \[
      \int_\Theta \varphi(\theta')\,S_t(\omega)(d\theta')
        = \mathbb E[\varphi(\theta)\mid\mathcal F_t](\omega)\qquad\text{a.s.}
    \]
    In particular:
    \textbf{(i)} $S_t$ is a sufficient statistic for $\theta$ given $\mathcal F_t$;
    \textbf{(ii)} since $\mathcal F_t = \sigma(Z_t)$, the Doob--Dynkin lemma implies that
    $S_t = f_t(Z_t)$ for some Borel measurable $f_t$;
    \textbf{(iii)} $\{S_t\}_{t\in\mathbb T}$ is a bounded $(\mathcal F_t)$-martingale.
\end{lemma}
\begin{proof}
    \uses{lem:existence, lem:sufficiency, lem:factorization, lem:martingale}
    命题是以下四步之合取：存在性与核性质（引理~\ref{lem:existence}）、充分性
    （引理~\ref{lem:sufficiency}，即 (i)）、由 $Z_t$ 因子分解（引理~\ref{lem:factorization}，即 (ii)）、
    鞅性质（引理~\ref{lem:martingale}，即 (iii)）。逐字证明见 \texttt{docs/paper-lemma1.md}。
\end{proof}

%======================================================================
% 步骤 1：存在性 + 核性质
%======================================================================
\begin{lemma}[Existence of the posterior kernel]\label{lem:existence}
    \lean{PosteriorProcess.kernel_exists}
    \uses{def:setup}
    对每个 $t$，存在 $\mathcal F_t$-可测的 $S_t:\Omega\to\Delta(\Theta)$，使得对每个有界 Borel
    $\varphi$，$\;\omega\mapsto\int_\Theta\varphi\,dS_t(\omega)$ 与
    $\mathbb E[\varphi(\theta)\mid\mathcal F_t]$ 几乎处处相等。
\end{lemma}
\begin{proof}
    $\Theta$ 紧 Polish $\Rightarrow$ 标准 Borel，由正则条件分布（Kallenberg 2002, Thm~6.3）即得。
    \emph{Lean 状态}：声明为 \texttt{PosteriorProcess.kernel\_exists}，**证明尚含 \texttt{sorry}**（未背书）。
\end{proof}

%======================================================================
% 步骤 2：(i) 充分性
%======================================================================
\begin{lemma}[Part (i): Sufficiency]\label{lem:sufficiency}
    \uses{lem:existence}
    在 $S_t$ 上取条件与在 $\mathcal F_t$ 上取条件给出 $\theta$ 相同的条件分布；即 $S_t$ 是
    $\theta$ 关于 $\mathcal F_t$ 的充分统计量。
\end{lemma}
\begin{proof}
    \uses{lem:footnote}
    对有界 Borel $\varphi$，$F(\mu):=\int\varphi\,d\mu$ 在 $\Delta(\Theta)$ 上 Borel 可测
    （泛函单调类定理：对连续 $\varphi$ 由弱拓扑定义成立，再由控制收敛延拓到有界 Borel）。
    故 $\mathbb E[\varphi(\theta)\mid\mathcal F_t]=F(S_t)$ 是 $\sigma(S_t)$-可测，由条件期望唯一性得两条件期望相等。
    \emph{Lean 状态}：**尚未形式化**（无对应声明）。\emph{另见}引理~\ref{lem:footnote}：论文此处一条脚注的错误已被 Lean 反证。
\end{proof}

%======================================================================
% 步骤 3：(ii) Doob–Dynkin —— 抽象层（已证）+ 库缺口 + 完整版
%======================================================================
\begin{lemma}[Abstract Doob--Dynkin factorization]\label{lem:doob-abstract}
    \lean{PosteriorProcess.factorization}\leanok
    若 $S$ 关于 $\sigma(Z)$ 可测、取值于标准 Borel 空间，则存在 Borel 可测 $f$ 使得 $S=f\circ Z$。
\end{lemma}
\begin{proof}\leanok
    Mathlib 的 \texttt{Measurable.exists\_eq\_measurable\_comp}。**已由 Lean 内核完整证明（无 \texttt{sorry}）。**
\end{proof}

\begin{lemma}[$\Delta(\Theta)$ is standard Borel]\label{lem:deltaTheta-sb}
    \lean{PosteriorProcess.deltaTheta_standardBorel}
    \uses{def:setup}
    $\Delta(\Theta)$（$\Theta$ 紧 Polish 上的概率测度，弱拓扑）是标准 Borel 空间。
\end{lemma}
\begin{proof}
    数学上成立（紧 Polish 上概率测度空间仍是紧 Polish）。\emph{Lean 状态}：当前 Mathlib
    缺 \texttt{ProbabilityMeasure} 的相应 \texttt{PolishSpace}/\texttt{StandardBorelSpace} 实例，
    故声明 \texttt{PosteriorProcess.deltaTheta\_standardBorel} 的**证明含 \texttt{sorry}（库缺口）**。
\end{proof}

\begin{lemma}[Part (ii): Induction by $Z_t$]\label{lem:factorization}
    \lean{PosteriorProcess.factorization_posterior}
    \uses{lem:existence, lem:doob-abstract, lem:deltaTheta-sb}
    由于 $\mathcal F_t=\sigma(Z_t)$，存在 Borel 可测 $f_t:E_t\to\Delta(\Theta)$ 使得
    $S_t=f_t(Z_t)$ a.s.；即公开历史 $Z_t$ 诱导后验 $S_t$。
\end{lemma}
\begin{proof}
    把抽象因子分解（引理~\ref{lem:doob-abstract}）用到 $\Delta(\Theta)$ 上，需后者为标准 Borel
    （引理~\ref{lem:deltaTheta-sb}）。\emph{Lean 状态}：声明为 \texttt{factorization\_posterior}，
    抽象步骤已证，但整体**经由引理~\ref{lem:deltaTheta-sb} 的库缺口仍含 \texttt{sorry}**。
\end{proof}

%======================================================================
% 步骤 4：(iii) 鞅性质（已证）
%======================================================================
\begin{lemma}[Part (iii): Martingale property]\label{lem:martingale}
    \lean{PosteriorProcess.testfun_martingale}\leanok
    \uses{lem:existence}
    对每个有界 Borel $\varphi$，实值过程 $t\mapsto\int_\Theta\varphi\,dS_t$ 是
    $(\mathcal F_t)$-鞅；从而 $\{S_t\}$ 是有界测度值鞅。
\end{lemma}
\begin{proof}\leanok
    由塔性质：$\mathbb E[\int\varphi\,dS_t\mid\mathcal F_s]
    =\mathbb E[\varphi(\theta)\mid\mathcal F_s]=\int\varphi\,dS_s$（$s\le t$）。Mathlib 的
    \texttt{martingale\_condExp}。**已由 Lean 内核完整证明（无 \texttt{sorry}）。**
\end{proof}

%======================================================================
% 审查：论文 (i) 脚注的错误 —— 已由 Lean 反证
%======================================================================
\begin{lemma}[Refutation of the Part (i) footnote]\label{lem:footnote}
    \lean{PosteriorProcess.indicatorRat_not_pointwiseLimit_continuous}\leanok
    有理数示性函数 $\mathbf 1_{\mathbb Q}$（有界 Borel）**不是**任何连续函数列的逐点极限。
    因此论文 (i) 的脚注断言“有界 Borel 函数 $=$ 一致有界连续函数列的逐点极限”**不成立**。
\end{lemma}
\begin{proof}\leanok
    自足的 Baire 纲反证：若 $g_n\to\mathbf 1_{\mathbb Q}$ 逐点，则 $\mathbb R$ 被可数个闭集
    $A_N=\{x:\forall n\ge N,\ \tfrac12\le g_n(x)\}$、$B_N=\{x:\forall n\ge N,\ g_n(x)\le\tfrac12\}$ 覆盖；
    由 Baire 纲某个含非空开集，其中的无理/有理点与极限值矛盾。**已由 Lean 内核完整证明（无 \texttt{sorry}）。**
    \emph{注}：这不影响 (i) 的结论——正文的泛函单调类论证本身正确（见引理~\ref{lem:sufficiency}）；
    该脚注只是一条方法有误且非必要的注解，建议在论文中删除。
\end{proof}
